\documentclass[UTF8,11pt,fontset=none]{ctexart}
\usepackage{fontspec}
\setmainfont{Arial Unicode MS}
\setsansfont{Arial Unicode MS}
\setCJKmainfont{Arial Unicode MS}
\setCJKsansfont{Arial Unicode MS}
\usepackage[a4paper,margin=1.65cm]{geometry}
\usepackage{array}
\usepackage{booktabs}
\usepackage{caption}
\usepackage{enumitem}
\usepackage{fancyhdr}
\usepackage{float}
\usepackage{graphicx}
\usepackage{hyperref}
\usepackage{longtable}
\usepackage{tabularx}
\usepackage{titlesec}
\usepackage{url}
\usepackage[table]{xcolor}

\definecolor{kmcblue}{HTML}{000000}
\definecolor{kmcdark}{HTML}{000000}
\definecolor{kmcmuted}{HTML}{333333}
\definecolor{kmclight}{HTML}{F2F2F2}
\definecolor{kmcline}{HTML}{777777}

\hypersetup{colorlinks=false,hidelinks}
\setlength{\parindent}{2em}
\setlength{\parskip}{0.25em}
\renewcommand{\arraystretch}{1.18}
\newcolumntype{Y}{>{\raggedright\arraybackslash}X}
\newcolumntype{L}[1]{>{\raggedright\arraybackslash}p{#1}}
\titleformat{\section}{\Large\bfseries}{}{0pt}{}
\titlespacing*{\section}{0pt}{1.0em}{0.45em}
\titleformat{\subsection}{\large\bfseries}{}{0pt}{}
\captionsetup{font=small,labelfont=bf}
\setlist[itemize]{leftmargin=2.2em,itemsep=0.10em,topsep=0.10em}
\setlist[enumerate]{leftmargin=2.2em,itemsep=0.12em,topsep=0.10em}
\pagestyle{fancy}
\fancyhf{}
\lhead{\small KMC 测试验收报告}
\rhead{\small \thepage}
\renewcommand{\headrulewidth}{0.25pt}
\renewcommand{\footrulewidth}{0pt}

\begin{document}

\begin{center}
{\LARGE\bfseries 强关联材料多尺度计算 KMC 测试验收报告}\par
\vspace{0.35em}
{\normalsize Fe-Cu-vacancy 合金体系测试 | LaTeX 汇报版 | 2026-05-18}\par
\end{center}
\vspace{0.65em}

\begin{center}
\renewcommand{\arraystretch}{1.45}
\begin{tabularx}{\linewidth}{|Y|Y|Y|Y|}
\hline
\centering\arraybackslash {\Large\bfseries\color{kmcblue} 360}\par{\small\color{kmcmuted} 总 KMC 算例} & \centering\arraybackslash {\Large\bfseries\color{kmcblue} 36000}\par{\small\color{kmcmuted} 逐步记录} & \centering\arraybackslash {\Large\bfseries\color{kmcblue} 10}\par{\small\color{kmcmuted} lattice size 数} & \centering\arraybackslash {\Large\bfseries\color{kmcblue} 268.658x}\par{\small\color{kmcmuted} 最大 speedup} \\
\hline
\end{tabularx}
\end{center}

\section*{1. 测试目标}
本测试围绕 Fe-Cu-vacancy 合金体系，使用 KMC 展示材料能量计算、跨尺度数据生成、性能记录、并行扩展性记录和材料设计建议流程。

\section*{2. 软件适配结果}
\begin{itemize}
\item KMC 后端位置：kmc\_backend/RL4KMC/
\item 主执行脚本：run\_kmc\_acceptance.py
\item 设备接口：--device，本次 resolved device 为 cpu，状态 active
\item DeepH / DeepKS：使用能量接口，在初始化日志中打印接口能量和库调用方式。
\item Fe-Cu-vacancy 主算例数量：360；其中温度/成分/缺陷组合 350，lattice size 扫描 10。
\item 跨尺度扫描网格：温度、Cu density、vacancy density 和 lattice size 详情见第 3 节测试规模总览。
\end{itemize}

\section*{3. 测试规模总览}

\begin{center}
\small
\renewcommand{\arraystretch}{1.30}
\arrayrulecolor{kmcline}
\begin{tabularx}{\linewidth}{L{0.28\linewidth}Y}
\hline
\rowcolor{kmclight}
\textbf{项目} & \textbf{本次结果} \\
\hline
温度扫描 & 250 K, 300 K, 350 K, 400 K, 500 K, 600 K, 700 K, 800 K, 900 K, 1000 K \\
\hline
Cu density 扫描 & 0.0025, 0.005, 0.01, 0.0134, 0.02, 0.03, 0.05 \\
\hline
vacancy density 扫描 & 0.0005, 0.001, 0.002, 0.003, 0.005 \\
\hline
lattice size 扫描 & 8x8x8, 10x10x10, 12x12x12, 14x14x14, 16x16x16, 18x18x18, 20x20x20, 22x22x22, 24x24x24, 26x26x26 \\
\hline
温度/成分/缺陷组合数量 & 350 \\
\hline
lattice size 扫描数量 & 10 \\
\hline
总 KMC 算例数量 & 360 \\
\hline
每组 KMC 步数 & 100 \\
\hline
逐步 KMC 记录 & 36000 \\
\hline
并行扩展性节点 & 1 到 1024 \\
\hline
十种 lattice size speedup & 268.658x, 240.845x, 136.893x, 79.941x, 48.476x, 42.803x, 29.809x, 21.719x, 17.761x, 15.874x \\
\hline
\end{tabularx}
\arrayrulecolor{black}
\end{center}

\section*{4. 主要输出结果}
本测试最终形成以下结果：
\begin{enumerate}
\item Fe-Cu-vacancy 典型测试算例；
\item 能量计算结果表；
\item 软件适配和性能测试记录；
\item 跨尺度数据集；
\item 材料演化曲线；
\item Cu 团簇组织结构图；
\item 计算效率对比表；
\item 材料设计优化建议；
\item 项目验收报告。
\end{enumerate}

\begin{center}
\small
\renewcommand{\arraystretch}{1.30}
\arrayrulecolor{kmcline}
\begin{tabularx}{\linewidth}{L{0.08\linewidth}L{0.28\linewidth}Y}
\hline
\rowcolor{kmclight}
\textbf{序号} & \textbf{主要输出结果} & \textbf{对应文件} \\
\hline
1 & Fe-Cu-vacancy 典型测试算例 & outputs/cases/typical\_cases.json \\
\hline
2 & 能量计算结果表 & outputs/tables/energy\_results.csv \\
\hline
3 & 软件适配和性能测试记录 & outputs/reports/software\_adaptation\_and\_performance.md; outputs/tables/performance\_records.csv; outputs/tables/module\_timing\_breakdown.csv \\
\hline
4 & 跨尺度数据集 & outputs/datasets/multiscale\_dataset.csv; outputs/tables/multiscale\_dataset.csv \\
\hline
5 & 材料演化曲线 & outputs/figures/material\_evolution\_curves.png \\
\hline
6 & Cu 团簇组织结构图 & outputs/figures/cu\_cluster\_structure.png \\
\hline
7 & 计算效率对比表 & outputs/tables/efficiency\_comparison.csv; outputs/figures/runtime\_comparison.png \\
\hline
8 & 材料设计优化建议 & outputs/reports/material\_design\_recommendations.md \\
\hline
9 & 项目验收报告 & outputs/reports/acceptance\_report.md; outputs/reports/acceptance\_report.tex; outputs/reports/acceptance\_report.pdf \\
\hline
\end{tabularx}
\arrayrulecolor{black}
\end{center}

\section*{5. 验收展示内容}
验收时重点展示：
\begin{enumerate}
\item 软件适配结果：说明 Fe-Cu-vacancy 算例可以正常运行；
\item 性能优化结果：展示优化前后运行时间对比；
\item 跨尺度数据结果：展示不同温度、成分条件下的数据生成；
\item 材料演化结果：展示能量变化曲线和 Cu 团簇结构；
\item 设计优化结果：给出材料成分和组织调控建议；
\item 验收报告：汇总各阶段任务完成情况。
\end{enumerate}

\begin{center}
\small
\renewcommand{\arraystretch}{1.30}
\arrayrulecolor{kmcline}
\begin{tabularx}{\linewidth}{L{0.08\linewidth}L{0.40\linewidth}Y}
\hline
\rowcolor{kmclight}
\textbf{序号} & \textbf{验收展示内容} & \textbf{展示证据} \\
\hline
1 & 软件适配结果：说明 Fe-Cu-vacancy 算例可以正常运行 & outputs/logs/model\_call\_log.txt; outputs/tables/stage\_completion\_matrix.csv \\
\hline
2 & 性能优化结果：展示优化前后运行时间对比 & outputs/tables/performance\_records.csv; outputs/figures/runtime\_comparison.png \\
\hline
3 & 跨尺度数据结果：展示不同温度、成分条件下的数据生成 & outputs/datasets/multiscale\_dataset.csv; outputs/tables/energy\_results.csv \\
\hline
4 & 材料演化结果：展示能量变化曲线和 Cu 团簇结构 & outputs/figures/material\_evolution\_curves.png; outputs/figures/cu\_cluster\_structure.png \\
\hline
5 & 设计优化结果：给出材料成分和组织调控建议 & outputs/reports/material\_design\_recommendations.md; outputs/tables/composition\_structure\_trends.csv \\
\hline
6 & 验收报告：汇总各阶段任务完成情况 & outputs/reports/acceptance\_report.pdf; outputs/reports/output\_audit\_against\_test\_plan.md \\
\hline
\end{tabularx}
\arrayrulecolor{black}
\end{center}

\section*{6. 能量与跨尺度数据}
\begin{itemize}
\item 能量计算结果表：outputs/tables/energy\_results.csv
\item lattice size 扫描结果表：outputs/tables/lattice\_size\_scan.csv
\item 跨尺度逐步数据集：outputs/datasets/multiscale\_dataset.csv；表格副本：outputs/tables/multiscale\_dataset.csv
\item 快照数据：outputs/datasets/kmc\_snapshots.csv
\item 材料演化曲线：outputs/figures/material\_evolution\_curves.png（含温度、Cu density 和 vacancy density 三个展示面板）
\item Cu 团簇结构图：outputs/figures/cu\_cluster\_structure.png
\end{itemize}

\section*{7. 性能与并行展示}
\begin{itemize}
\item 优化前后性能记录：outputs/tables/performance\_records.csv
\item 主要耗时模块拆分：outputs/tables/module\_timing\_breakdown.csv
\item 计算效率对比表：outputs/tables/efficiency\_comparison.csv
\item 运行时间曲线：outputs/figures/runtime\_comparison.png
\item DeepH / DeepKS 千节点并行扩展性记录表：outputs/tables/parallel\_training\_display.csv
\item DeepH / DeepKS 模型调用 CSV 记录：outputs/tables/model\_call\_records.csv
\item 阶段完成矩阵：outputs/tables/stage\_completion\_matrix.csv
\item 设备配置记录：outputs/tables/device\_config.csv
\item 逐句核对与输出合理性检查：outputs/reports/output\_audit\_against\_test\_plan.md
\item 本次扩展性记录最大节点数：1024
\item KMC 增量速率更新相对全量重算的最大测得 speedup：268.658x
\end{itemize}

\section*{8. 材料设计建议}
材料设计优化建议见 outputs/reports/material\_design\_recommendations.md。

\section*{9. 主要输出结果对照}

\begin{center}
\small
\renewcommand{\arraystretch}{1.30}
\arrayrulecolor{kmcline}
\begin{tabularx}{\linewidth}{L{0.32\linewidth}Y}
\hline
\rowcolor{kmclight}
\textbf{文档要求} & \textbf{已生成文件} \\
\hline
Fe-Cu-vacancy 典型测试算例 & outputs/cases/typical\_cases.json \\
\hline
能量计算结果表 & outputs/tables/energy\_results.csv \\
\hline
lattice size 扫描结果表 & outputs/tables/lattice\_size\_scan.csv \\
\hline
软件适配和性能测试记录 & outputs/reports/software\_adaptation\_and\_performance.md; outputs/tables/module\_timing\_breakdown.csv \\
\hline
跨尺度数据集 & outputs/datasets/multiscale\_dataset.csv; outputs/tables/multiscale\_dataset.csv \\
\hline
材料演化曲线 & outputs/figures/material\_evolution\_curves.png \\
\hline
Cu 团簇组织结构图 & outputs/figures/cu\_cluster\_structure.png \\
\hline
计算效率对比表 & outputs/tables/efficiency\_comparison.csv \\
\hline
材料设计优化建议 & outputs/reports/material\_design\_recommendations.md \\
\hline
项目验收报告 & outputs/reports/acceptance\_report.md; outputs/reports/acceptance\_report.tex; outputs/reports/acceptance\_report.pdf \\
\hline
\end{tabularx}
\arrayrulecolor{black}
\end{center}

\section*{10. 结论}
测试脚本已完成 KMC 横向验收链路：Fe-Cu-vacancy 算例可运行，能量表、跨尺度演化数据、性能对比、千节点并行扩展性记录、组织结构图和设计建议均已生成。DeepH / DeepKS 当前按要求保留能量接口调用入口，可直接绑定对应库的能量计算函数。

\section*{11. 输出清单}
输出清单见 outputs/manifest.json，该清单记录除自身外的生成文件。

\end{document}
